\documentclass[11pt,a4paper]{article}

\usepackage[T1]{fontenc}
\usepackage[utf8]{inputenc}
\usepackage[brazilian]{babel}
\usepackage[margin=2.2cm]{geometry}
\usepackage{xcolor}
\usepackage{tcolorbox}
\tcbuselibrary{skins,breakable}
\usepackage{tikz}
\usetikzlibrary{arrows.meta,positioning,fit,backgrounds,shapes.multipart}
\usepackage{listings}
\usepackage{hyperref}
\usepackage{enumitem}
\usepackage{fancyhdr}
\usepackage{booktabs}
\usepackage{longtable}
\usepackage{parskip}
\usepackage{titlesec}

% ── Paleta ──────────────────────────────────────────────────────────────────
\definecolor{droneblue}{HTML}{2563EB}
\definecolor{teamgreen}{HTML}{15803D}
\definecolor{alertred}{HTML}{DC2626}
\definecolor{ackpurple}{HTML}{7C3AED}
\definecolor{panelbg}{HTML}{F3F4F6}
\definecolor{codebg}{HTML}{1E1E2E}
\definecolor{codefg}{HTML}{CDD6F4}
\definecolor{codegreen}{HTML}{A6E3A1}
\definecolor{codeyellow}{HTML}{F9E2AF}

\titleformat{\section}{\Large\bfseries\color{droneblue}}{\thesection}{0.5em}{}
\titleformat{\subsection}{\large\bfseries\color{teamgreen!70!black}}{\thesubsection}{0.5em}{}

\hypersetup{colorlinks=true, linkcolor=droneblue, urlcolor=droneblue, citecolor=droneblue,
  pdftitle={ECHOSAR-Net --- Guia Ilustrado dos Logs de Simulação},
  pdfauthor={Rodrigo Paczkovski}}

\lstdefinestyle{simlog}{
  backgroundcolor=\color{codebg},
  basicstyle=\ttfamily\scriptsize\color{codefg},
  breaklines=true,
  columns=fullflexible,
  keepspaces=true,
  frame=single,
  framerule=0pt,
  rulecolor=\color{codebg},
  xleftmargin=4pt, xrightmargin=4pt,
  aboveskip=6pt, belowskip=6pt,
  literate=%
    {á}{{\'a}}1 {é}{{\'e}}1 {í}{{\'i}}1 {ó}{{\'o}}1 {ú}{{\'u}}1
    {ã}{{\~a}}1 {õ}{{\~o}}1 {â}{{\^a}}1 {ê}{{\^e}}1 {ô}{{\^o}}1
    {ç}{{\c c}}1 {→}{{$\rightarrow$}}1
    {Á}{{\'A}}1 {É}{{\'E}}1 {Í}{{\'I}}1 {Ó}{{\'O}}1 {Ú}{{\'U}}1
    {Ã}{{\~A}}1 {Õ}{{\~O}}1 {Â}{{\^A}}1 {Ê}{{\^E}}1 {Ô}{{\^O}}1
    {Ç}{{\c C}}1,
}
\lstdefinestyle{shellcmd}{
  backgroundcolor=\color{panelbg},
  basicstyle=\ttfamily\small,
  breaklines=true,
  columns=fullflexible,
  frame=single,
  framerule=0.4pt,
  rulecolor=\color{black!30},
  xleftmargin=4pt, xrightmargin=4pt,
  aboveskip=6pt, belowskip=6pt,
}

\newtcolorbox{explica}[1][]{
  colback=panelbg, colframe=droneblue!60, boxrule=0.6pt,
  arc=2pt, left=6pt, right=6pt, top=4pt, bottom=4pt,
  fonttitle=\bfseries\small, title={#1}, breakable
}

\newcommand{\droneTag}{\textcolor{droneblue}{\textbf{[DRONE]}}}
\newcommand{\teamTag}{\textcolor{teamgreen}{\textbf{[TEAM]}}}

\pagestyle{fancy}
\fancyhf{}
\fancyhead[L]{\small ECHOSAR-Net --- Guia dos Logs}
\fancyhead[R]{\small \thepage}
\renewcommand{\headrulewidth}{0.4pt}

\begin{document}

% ══════════════════════════════════════════════════════════════════════════
\begin{titlepage}
\centering
\vspace*{2cm}
{\Huge\bfseries\color{droneblue} ECHOSAR-Net}\\[0.4cm]
{\Large\color{teamgreen!70!black} Guia Ilustrado dos Logs de Simulação}\\[1.2cm]

\begin{tikzpicture}[
  drone/.style={circle, fill=droneblue, minimum size=10pt, inner sep=0pt},
  team/.style={rectangle, fill=teamgreen, minimum size=10pt, inner sep=0pt, rounded corners=1pt},
]
  \draw[droneblue!20, fill=droneblue!5, rounded corners=8pt] (-5.5,-3) rectangle (5.5,3);
  \node[drone] (d1) at (-3.5,2) {};
  \node[drone] (d2) at (-1,2.3) {};
  \node[drone] (d3) at (1.5,1.8) {};
  \node[drone] (d4) at (3.8,2.2) {};
  \node[drone] (d5) at (0,0.8) {};
  \node[team]  (t1) at (-3,-1.8) {};
  \node[team]  (t2) at (0.5,-2.2) {};
  \node[team]  (t3) at (3.5,-1.5) {};

  \draw[droneblue!50, thin] (d1) -- (d2);
  \draw[droneblue!50, thin] (d2) -- (d3);
  \draw[droneblue!50, thin] (d3) -- (d4);
  \draw[droneblue!50, thin] (d2) -- (d5);
  \draw[droneblue!50, thin] (d3) -- (d5);

  \draw[alertred, -{Latex[length=2mm]}, thick] (d5) -- (t2)
    node[midway, above, sloped, font=\tiny\bfseries] {VictimAlert};
  \draw[ackpurple, -{Latex[length=2mm]}, thick, dashed] (t2) to[bend left=15] node[midway, below, sloped, font=\tiny\bfseries] {VictimAck} (d5);
  \draw[teamgreen, -{Latex[length=2mm]}, thin] (t1) -- (d1) node[midway, above, sloped, font=\tiny] {TeamUpdate};
\end{tikzpicture}

\vspace{0.8cm}
{\large Como ler, entender e explicar cada linha de log da simulação}\\[0.3cm]
{\small Simulação OMNeT++ 6.2.0 + INET 4.5.4 --- Dissertação PPGCAP/UDESC}

\vfill
{\small Rodrigo Paczkovski \\ \today}
\end{titlepage}

\tableofcontents
\newpage

% ══════════════════════════════════════════════════════════════════════════
\section{Por que este guia existe}

Este documento explica, item por item, o que aparece nos logs quando a simulação
ECHOSAR-Net roda. O objetivo é permitir apresentar e defender cada linha de log
para a banca/orientação: o que ela significa, de onde vem no código, e por que
aquele valor apareceu ali.

\begin{explica}[Ideia central da simulação]
Drones sobrevoam uma área simulando busca e resgate (SAR). Não há vítimas físicas
nem sensores reais de detecção --- os drones geram \textbf{alertas sintéticos}
em intervalos aleatórios para gerar tráfego de rede representativo. O que se mede
é a \textbf{qualidade da comunicação} entre drones e equipes de resgate em terra:
o alerta chega? Em quanto tempo? Quantas retransmissões precisou?
\end{explica}

\subsection{Como gerar um log legível}

Por padrão, a simulação roda em modo ``express'' (rápido, sem imprimir eventos)
para não desperdiçar tempo de CPU. Para ver as mensagens de aplicação
(\texttt{[DRONE]}, \texttt{[TEAM]}), é preciso pedir o nível \texttt{INFO}
explicitamente:

\begin{lstlisting}[style=shellcmd]
./run.sh --info -c SmokeTest_Beacons
\end{lstlisting}

\begin{explica}[O que a flag \texttt{--info} faz de verdade]
O script \texttt{run.sh} passa duas flags ao motor de simulação: \texttt{--cmdenv-log-level}
(qual severidade mostrar) e \texttt{--cmdenv-express-mode} (se os eventos são
impressos ou não). \textbf{As duas precisam estar certas} --- só pedir
\texttt{log-level=INFO} não basta se o express-mode estiver ligado, porque ele
suprime a impressão de qualquer evento independente do nível. \texttt{--info}
cuida das duas coisas.
\end{explica}

\newpage
% ══════════════════════════════════════════════════════════════════════════
\section{Quem fala com quem: a arquitetura de mensagens}

Cinco tipos de mensagem UDP circulam na rede, cada uma numa porta fixa
(\texttt{src/app/ports.h}):

\begin{center}
\begin{tabular}{@{}lllp{6.2cm}@{}}
\toprule
\textbf{Porta} & \textbf{Mensagem} & \textbf{Sentido} & \textbf{Quando é enviada} \\
\midrule
5000 & \texttt{VictimAlert} & drone $\to$ equipe & Drone gera alerta sintético e escolhe a equipe mais próxima conhecida \\
5001 & \texttt{TeamUpdate}  & equipe $\to$ broadcast & Beacon periódico da equipe (posição + IP) \\
5002 & \texttt{VictimAck}   & equipe $\to$ drone origem & Equipe confirma recebimento do alerta \\
5003 & \texttt{DroneStatus} & drone $\to$ equipe & Resposta do drone ao receber um \texttt{TeamUpdate} \\
5004 & \texttt{VictimAlert} & drone $\to$ drone (relay) & Drone repassa alerta em broadcast quando não conhece nenhuma equipe \\
\bottomrule
\end{tabular}
\end{center}

\subsection{Diagrama de sequência do fluxo completo}

\begin{center}
\begin{tikzpicture}[
  actor/.style={font=\bfseries\small},
  lifeline/.style={draw, thick, -},
  msg/.style={-{Latex[length=2.5mm]}, thick},
  every node/.style={font=\small},
]
  \node[actor, teamgreen] (T) at (0,0) {EQUIPE};
  \node[actor, droneblue] (D) at (7,0) {DRONE};
  \draw[lifeline, teamgreen!40] (T) -- ++(0,-9.3);
  \draw[lifeline, droneblue!40] (D) -- ++(0,-9.3);

  \draw[msg, teamgreen] (0,-1) -- node[above, sloped]{\scriptsize\texttt{TeamUpdate} bcast :5001} (7,-1.5);
  \draw[msg, droneblue] (7,-2.3) -- node[above, sloped]{\scriptsize\texttt{DroneStatus} uni :5003} (0,-2.8);

  \node[alertred, font=\scriptsize\itshape] at (3.5,-3.6) {drone gera alerta sintético (Exp(alertInterval))};
  \draw[msg, alertred] (7,-4.1) -- node[above, sloped]{\scriptsize\texttt{VictimAlert} uni :5000} (0,-4.6);
  \draw[msg, ackpurple, dashed] (0,-5.4) -- node[above, sloped]{\scriptsize\texttt{VictimAck} uni :5002} (7,-5.9);

  \node[font=\scriptsize\itshape] at (3.5,-6.7) {sem ACK em \texttt{retryInterval}=10s $\to$ retry p/ outra equipe (até \texttt{maxRetries}=5)};

  \node[droneblue, font=\scriptsize\itshape] at (7,-7.6) {[se nenhuma equipe conhecida]};
  \node[actor, droneblue] (V) at (9.6,-8.3) {vizinhos};
  \draw[lifeline, droneblue!40] (V) -- ++(0,-0.9);
  \draw[msg, droneblue] (7,-8.6) -- node[above, sloped]{\scriptsize\texttt{VictimAlert} bcast :5004} (V);
\end{tikzpicture}
\end{center}

\begin{explica}[Por que o ACK volta direto ao drone que criou o alerta, não a quem repassou]
O campo \texttt{originIp} dentro do \texttt{VictimAlert} guarda o IP do drone que
\textbf{originou} o alerta. Mesmo que a mensagem tenha passado por vários relays
até chegar à equipe, o \texttt{VictimAck} é endereçado direto a esse IP de
origem via roteamento AODV multi-hop --- não ao último drone que fez relay.
\end{explica}

\newpage
% ══════════════════════════════════════════════════════════════════════════
\section{Anatomia de cada mensagem}

\subsection{\droneTag~TeamUpdate --- beacon da equipe}

\begin{lstlisting}[style=simlog]
[INFO]	[TEAM team0] TeamUpdate broadcast (ip=10.0.0.2)
[INFO]	[DRONE drone[0]] tabela: team0 ip=10.0.0.2
\end{lstlisting}

\begin{explica}[Campos do \texttt{TeamUpdateChunk} --- 128 bytes]
\texttt{teamId} (nome da equipe) · \texttt{ipAddress} (IP dela na sub-rede
\texttt{10.0.0.x/24}) · \texttt{posX}/\texttt{posY} (posição cartesiana no
playground, em metros). \\[2pt]
\textbf{O que acontece:} a equipe transmite isso em broadcast a cada
\texttt{sendInterval} segundos (com um jitter aleatório inicial para não
colidir com outras equipes no rádio). Todo drone que recebe atualiza sua
\texttt{teamTable} --- é assim que o drone \emph{descobre} quais equipes
existem e onde estão, sem nenhuma configuração estática.
\end{explica}

\subsection{\droneTag~DroneStatus --- resposta do drone}

\begin{lstlisting}[style=simlog]
[INFO]	[TEAM team0] DroneStatus de drone[0] pos=(2500,2500,100) RTT=0.004653659192s
\end{lstlisting}

\begin{explica}[Campos do \texttt{DroneStatusChunk} --- 128 bytes]
\texttt{droneId} · \texttt{posX}/\texttt{posY}/\texttt{posZ} (posição 3D,
incluindo altitude) · \texttt{sentAt} (timestamp de envio, usado para calcular
o RTT mostrado no log). \\[2pt]
\textbf{O que acontece:} é a resposta automática do drone a todo
\texttt{TeamUpdate} recebido --- funciona como um ``ping'' de posição. A
equipe usa isso para saber que drones estão por perto e ativos (não é
usado para nenhuma métrica de entrega, é diagnóstico de conectividade).
\end{explica}

\subsection{\droneTag~VictimAlert --- o alerta}

\begin{lstlisting}[style=simlog]
[INFO]	[DRONE drone[0]] alerta sintético gerado → drone[0]_1
[INFO]	[DRONE drone[0]] VictimAlert drone[0]_1 → team0 (10.0.0.2)
[INFO]	[TEAM team0] *** ALERTA de drone[0] msgId=drone[0]_1 vitima em (2500,2500) delay=0.000134368468s
\end{lstlisting}

\begin{explica}[Campos do \texttt{VictimAlertChunk} --- 256 bytes]
\texttt{droneId} · \texttt{msgId} (identificador único \texttt{droneId\_N},
usado para deduplicação) · \texttt{originIp} (IP de quem originou, para o ACK
voltar certo) · \texttt{posX}/\texttt{posY} (posição da ``vítima'') ·
\texttt{sentAt}. \\[2pt]
\textbf{O que acontece:} a cada intervalo exponencial (\texttt{alertInterval}),
o drone gera um alerta sintético e escolhe a equipe conhecida \textbf{mais
próxima} da sua \texttt{teamTable} para enviar via unicast. Se não conhece
nenhuma equipe, faz relay em broadcast (porta 5004) para drones vizinhos
tentarem entregar.
\end{explica}

\subsection{\teamTag~VictimAck --- a confirmação}

\begin{lstlisting}[style=simlog]
[INFO]	[TEAM team0] VictimAck drone[0]_1 → drone origem 10.0.0.1
[INFO]	[DRONE drone[0]] VictimAck recebido para drone[0]_1 de team0
\end{lstlisting}

\begin{explica}[Campos do \texttt{VictimAckChunk} --- 64 bytes]
\texttt{msgId} (referencia o alerta confirmado) · \texttt{teamId} (quem
confirmou) · \texttt{sentAt}. \\[2pt]
\textbf{O que acontece:} ao receber e processar um \texttt{VictimAlert}, a
equipe manda esse ACK de volta ao \texttt{originIp}. O drone remove o alerta
de \texttt{pendingAlerts} e para de retransmiti-lo. Se o mesmo \texttt{msgId}
chegar de novo (ACK perdido no caminho), a equipe reenvia o ACK sem contar de
novo nas métricas --- evita retransmitir o alerta inteiro por perda só do ACK.
\end{explica}

\newpage
% ══════════════════════════════════════════════════════════════════════════
\section{Por baixo do capô: rádio, MAC e roteamento AODV}

Tudo que foi mostrado até aqui (\texttt{[DRONE]}/\texttt{[TEAM]}) é log de
\textbf{aplicação}. Mas cada uma dessas mensagens desce por várias camadas de
rede simuladas antes de virar sinal de rádio --- e o \texttt{--info} imprime
o que acontece em \emph{todas} elas. Esta seção explica as linhas mais
frequentes desses níveis mais baixos.

\subsection{Camada física (PHY) --- rádio 802.11}

\begin{lstlisting}[style=simlog]
[INFO]	Radio mode changed from TRANSMITTER to RECEIVER.
[INFO]	Changing radio reception state from UNDEFINED to IDLE.
[INFO]	Changing radio transmission state from IDLE to UNDEFINED.
[INFO]	bpsk snr=6610.06 ber=0
[INFO]	16-Qam snr=6610.06 ber=6.97908e-290
[INFO]	Reception ended: successfully for (...)WlanAck (34 us 27 B) ...
	power = 661.006 pW, centerFrequency = 2.412 GHz, bandwidth = 20 MHz
[INFO]	Sending up (inet::Packet)WlanAck (14 B)
[INFO]	Changing radio reception state from RECEIVING to IDLE.
\end{lstlisting}

\begin{explica}[O que cada linha quer dizer]
\textbf{Radio mode / reception state / transmission state}: máquina de
estados do rádio (\texttt{IDLE}, \texttt{RECEIVING}, \texttt{TRANSMITTER}...)
--- um rádio half-duplex só pode estar em um estado por vez. \\
\textbf{\texttt{bpsk}/\texttt{16-Qam snr=... ber=...}}: para \emph{cada}
modulação suportada, o simulador calcula a \textbf{SNR} (relação sinal-ruído,
linear) e a \textbf{BER} (taxa de erro de bit) que aquele link teria se
usasse essa modulação --- serve para decidir se o quadro é decodificado com
sucesso. \\
\textbf{power = ... pW}: potência recebida do sinal, em picowatts --- compare
com a sensibilidade do receptor (\texttt{-85 dBm} no cenário) para saber se o
quadro sequer é detectável.
\end{explica}

\subsection{Camada de enlace (MAC) --- handshake 802.11}

\begin{lstlisting}[style=simlog]
[INFO]	Received frame from PHY: (inet::Packet)WlanAck (14 B)
[INFO]	Processing lower frame: WlanAck
[INFO]	Finishing last frame sequence step: history = ({} LAST-FRAME ACK)
[INFO]	Processing received frame WlanAck as originator in frame sequence.
[INFO]	Frame sequence finished.
[INFO]	Channel released.
\end{lstlisting}

\begin{explica}[Por que todo pacote de aplicação vira um \texttt{WlanAck}]
O IEEE 802.11 exige confirmação em nível de enlace para \textbf{todo} quadro
unicast: quem envia um \texttt{DATA} espera um \texttt{ACK} do MAC antes de
liberar o canal (\texttt{Channel released}). Isso é diferente do
\texttt{VictimAck} da aplicação --- este ACK é interno ao rádio, acontece em
microssegundos, e existe até para um simples \texttt{TeamUpdate}. Se o ACK
802.11 não chegar, o MAC retransmite no nível de enlace, \emph{antes} de
qualquer retry da aplicação entrar em ação.
\end{explica}

\subsection{Roteamento AODV --- descoberta de rota}

Quando um drone precisa mandar algo para um IP com o qual ainda não tem rota
ativa, o AODV descobre a rota sob demanda. Exemplo real de
\texttt{10.0.0.1} (drone) descobrindo rota para \texttt{10.0.0.2} (equipe)
para poder enviar um \texttt{DroneStatus}:

\begin{lstlisting}[style=simlog]
[INFO]	Finding route for source <unspec> with destination 10.0.0.2
[INFO]	Missing route for destination 10.0.0.2
[INFO]	Starting route discovery with originator 10.0.0.1 and destination 10.0.0.2
[INFO]	Sending a Route Request with target 10.0.0.2 and TTL= 2
	...
[INFO]	AODV Route Request arrived with source addr: 10.0.0.1 originator addr: 10.0.0.1 destination addr: 10.0.0.2
[INFO]	I am the destination node for which the route was requested
[INFO]	Sending Route Reply to 10.0.0.1
	...
[INFO]	AODV Route Reply arrived with source addr: 10.0.0.2 originator addr: 10.0.0.1 destination addr: 10.0.0.2
[INFO]	The Route Reply has arrived for our Route Request to node 10.0.0.2
[INFO]	Routing (inet::Packet)DroneStatus (156 B) with destination = 10.0.0.2,
	output interface = wlan0, next hop address = 10.0.0.2
\end{lstlisting}

\begin{explica}[Leitura linha a linha]
\begin{enumerate}[leftmargin=1.3em, itemsep=2pt]
  \item \texttt{Missing route} --- o drone não tem entrada na tabela de
  roteamento para \texttt{10.0.0.2}, então dispara descoberta.
  \item \textbf{RREQ} (\emph{Route Request}) sai em broadcast com
  \texttt{TTL=2} --- alcança até 2 saltos.
  \item Quem recebe o RREQ e \textbf{é} o destino responde direto com
  \textbf{RREP} (\emph{Route Reply}); quem não é, encaminharia o RREQ adiante
  (aqui não precisou, são vizinhos diretos).
  \item Ao receber o RREP, o drone grava a rota (\texttt{add route}) e agora
  pode \textbf{rotear} o \texttt{DroneStatus} de verdade, com \texttt{next
  hop} definido.
\end{enumerate}
Isso só acontece \textbf{uma vez} por par origem/destino (enquanto a rota
estiver ativa --- \texttt{activeRouteTimeout=5s}). Depois disso, os pacotes
seguintes só aparecem como \texttt{Active route found}, sem todo o handshake
RREQ/RREP de novo.
\end{explica}

\newpage
% ══════════════════════════════════════════════════════════════════════════
\section{Timers e parâmetros de controle}

\begin{center}
\begin{tabular}{@{}llp{7.2cm}@{}}
\toprule
\textbf{Parâmetro} & \textbf{Onde} & \textbf{Papel} \\
\midrule
\texttt{alertInterval} & drone (.ned) & Média da distribuição exponencial entre alertas sintéticos gerados \\
\texttt{teamTimeout}   & drone (.ned) & Se uma equipe some da \texttt{teamTable} por esse tempo sem novo \texttt{TeamUpdate}, é removida \\
\texttt{retryInterval} & drone (.ned) & Intervalo entre tentativas de reenvio de um alerta sem ACK \\
\texttt{maxRetries}    & drone (.ned) & Máximo de retransmissões antes do alerta ser descartado (\texttt{alertsExpired}) \\
\texttt{sendInterval}  & equipe (.ned) & Intervalo entre beacons \texttt{TeamUpdate} \\
\texttt{beaconJitter}  & equipe (.ned) & Atraso aleatório do primeiro beacon --- evita colisão MAC entre equipes \\
\bottomrule
\end{tabular}
\end{center}

\begin{explica}[Por que o \texttt{beaconJitter} é crítico]
Se todas as equipes transmitirem o primeiro beacon exatamente no mesmo
instante, o rádio 802.11 sofre colisões simultâneas e \textbf{nenhum} drone
recebe \texttt{TeamUpdate} --- a \texttt{teamTable} fica vazia e o PDR
despenca para perto de zero. O jitter espalha as transmissões iniciais no
tempo. O código rejeita (\texttt{cRuntimeError}) se \texttt{beaconJitter=0}
exatamente por esse motivo já ter causado uma regressão real na simulação.
\end{explica}

\subsection{Timers internos do drone (\texttt{SimpleDroneApp})}

\begin{itemize}[leftmargin=1.4em]
  \item \textbf{\texttt{alertTimer}} --- dispara \texttt{generateAlert()} a cada Exp(\texttt{alertInterval})
  \item \textbf{\texttt{timeoutTimer}} --- dispara \texttt{checkTimeouts()} a cada \texttt{teamTimeout}, limpando equipes inativas
  \item \textbf{\texttt{retryTimer}} --- dispara \texttt{retryPending()} a cada \texttt{retryInterval}, reenviando alertas sem ACK
\end{itemize}

\newpage
% ══════════════════════════════════════════════════════════════════════════
\section{Log real anotado --- uma simulação inteira em miniatura}

Trecho real de \texttt{SmokeTest\_Beacons} (1 drone parado + 2 equipes,
$t=0$ a $15$s) rodado com \texttt{./run.sh --info -c SmokeTest\_Beacons}.
Cada bloco abaixo é um ciclo completo do protocolo:

\begin{lstlisting}[style=simlog]
[INFO]	[TEAM team0] TeamUpdate broadcast (ip=10.0.0.2)
[INFO]	[DRONE drone[0]] tabela: team0 ip=10.0.0.2
[INFO]	[TEAM team0] DroneStatus de drone[0] pos=(2500,2500,100) RTT=0.004653659192s
[INFO]	[TEAM team1] TeamUpdate broadcast (ip=10.0.0.3)
[INFO]	[DRONE drone[0]] tabela: team1 ip=10.0.0.3
[INFO]	[TEAM team1] DroneStatus de drone[0] pos=(2500,2500,100) RTT=0.001449209234s
[INFO]	[DRONE drone[0]] alerta sintético gerado → drone[0]_1
[INFO]	[DRONE drone[0]] VictimAlert drone[0]_1 → team0 (10.0.0.2)
[INFO]	[TEAM team0] *** ALERTA de drone[0] msgId=drone[0]_1 vitima em (2500,2500) delay=0.000134368468s
[INFO]	[TEAM team0] VictimAck drone[0]_1 → drone origem 10.0.0.1
[INFO]	[DRONE drone[0]] VictimAck recebido para drone[0]_1 de team0
\end{lstlisting}

\begin{explica}[Leitura linha a linha]
\begin{enumerate}[leftmargin=1.3em, itemsep=2pt]
  \item As duas equipes anunciam sua posição/IP (\texttt{TeamUpdate}) em
  momentos diferentes --- efeito do \texttt{beaconJitter}.
  \item O drone confirma a recepção de cada beacon com \texttt{DroneStatus}
  --- por isso aparecem dois RTTs, um por equipe.
  \item O drone gera um alerta sintético (\texttt{drone[0]\_1}) e escolhe
  \texttt{team0} --- nesta simulação, a única candidata elegível no momento
  (equipe mais próxima entre as conhecidas).
  \item A equipe recebe, imprime \texttt{*** ALERTA ***} (marcador visual de
  entrega bem-sucedida) e calcula o \texttt{delay} --- tempo entre
  \texttt{sentAt} do alerta e o instante de recepção. Aqui, $0{,}13$~ms:
  drone e equipe estão a poucos metros um do outro no cenário de teste.
  \item A equipe devolve o \texttt{VictimAck} ao IP de origem do drone.
  \item O drone recebe o ACK e remove o alerta de \texttt{pendingAlerts} ---
  ciclo fechado, sem retransmissão necessária.
\end{enumerate}
\end{explica}

\newpage
% ══════════════════════════════════════════════════════════════════════════
\section{Métricas gravadas ao final (\texttt{finish()})}

Ao fim de cada simulação, cada módulo grava contadores (\texttt{scalar}) no
arquivo \texttt{.sca} --- é isso que \texttt{analysis/process\_results.py}
lê para calcular as métricas da dissertação.

\subsection{Do lado do drone}

\begin{center}
\begin{tabular}{@{}llp{6.7cm}@{}}
\toprule
\textbf{Scalar} & \textbf{Tipo} & \textbf{Significado} \\
\midrule
\texttt{alertsGenerated}  & contador & alertas sintéticos criados \\
\texttt{alertsSentDirect} & contador & \texttt{VictimAlert} enviados direto a uma equipe \\
\texttt{alertsSentRelay}  & contador & \texttt{VictimAlert} enviados em relay broadcast \\
\texttt{alertsRelayed}    & contador & alertas de \emph{outro} drone recebidos e repassados \\
\texttt{alertsAcked}      & contador & \texttt{VictimAck} recebidos (ciclo completo confirmado) \\
\texttt{alertsExpired}    & contador & descartados após \texttt{maxRetries} sem ACK \\
\texttt{totalRetries}     & contador & soma de tentativas de store-forward \\
\texttt{meanRTT}          & média & tempo alerta$\to$ACK, do lado do drone (ciclo completo, \textbf{não} é o atraso 1-via) \\
\bottomrule
\end{tabular}
\end{center}

\subsection{Do lado da equipe}

\begin{center}
\begin{tabular}{@{}llp{6.7cm}@{}}
\toprule
\textbf{Scalar} & \textbf{Tipo} & \textbf{Significado} \\
\midrule
\texttt{alertsReceived}      & contador & \texttt{VictimAlert} únicos recebidos (após deduplicação) \\
\texttt{teamUpdatesSent}     & contador & beacons enviados \\
\texttt{droneStatusReceived} & contador & respostas de drones recebidas \\
\texttt{meanDeliveryDelay}   & média & atraso de entrega \textbf{1-via} (drone$\to$equipe) --- esta é a métrica de ``atraso fim-a-fim'' reportada na dissertação \\
\bottomrule
\end{tabular}
\end{center}

\begin{explica}[PDR vs. AppACK --- por que existem duas métricas de sucesso]
\textbf{PDR} = \texttt{alertsReceived}/\texttt{alertsGenerated} mede
\emph{chegada ao destino} (a equipe recebeu). \textbf{AppACK} =
\texttt{alertsAcked}/\texttt{alertsGenerated} mede o \emph{ciclo completo}
confirmado (o drone soube que chegou). Com AODV multi-hop os dois quase
coincidem, porque o \texttt{VictimAck} também consegue voltar vários saltos
--- a diferença residual vem de ACKs perdidos no caminho de volta.
\end{explica}

\newpage
% ══════════════════════════════════════════════════════════════════════════
\section{Glossário rápido}

\begin{longtable}{@{}lp{10.5cm}@{}}
\toprule
\textbf{Termo} & \textbf{Significado} \\
\midrule
\endhead
AODV & Protocolo de roteamento reativo multi-hop --- descobre rotas sob demanda (RREQ/RREP), sem tabelas estáticas \\
RREQ / RREP & \emph{Route Request} / \emph{Route Reply} --- mensagens do AODV para descobrir e confirmar uma rota \\
SNR / BER & Relação sinal-ruído / taxa de erro de bit --- determinam se um quadro de rádio é decodificado com sucesso \\
ARP & Protocolo que resolve IP $\to$ endereço MAC antes de enviar um quadro no enlace \\
FANET & Flying Ad-hoc Network --- rede ad hoc formada pelos próprios drones em voo \\
Relay & Repasse de um alerta por um drone que não tem equipe conhecida, esperando que um vizinho tenha \\
Store-forward & Estratégia de reter o alerta e reenviar periodicamente até receber ACK ou esgotar tentativas \\
Dedup / \texttt{seenAlerts} & Conjunto de \texttt{msgId} já vistos, evita processar/retransmitir o mesmo alerta duas vezes \\
Express-mode & Modo do Cmdenv que roda sem imprimir eventos (rápido); precisa ser desligado para ver logs \\
EV\_INFO / EV\_WARN & Macros do OMNeT++ para logging em diferentes níveis de severidade \\
Scalar (\texttt{.sca}) & Arquivo com os contadores finais de cada módulo, um por seed \\
Seed / run & Uma execução independente da simulação com número aleatório diferente --- várias seeds dão média $\pm$ desvio \\
\bottomrule
\end{longtable}

\vspace{1cm}
\begin{center}
\small\itshape
Gerado a partir do código em \texttt{src/app/\{SimpleDroneApp,SimpleTeamApp\}.\{h,cc,ned\}} \\
e de uma execução real de \texttt{SmokeTest\_Beacons} --- ver \texttt{CLAUDE.md} para detalhes de arquitetura.
\end{center}

\end{document}
